% Options for packages loaded elsewhere
\PassOptionsToPackage{unicode}{hyperref}
\PassOptionsToPackage{hyphens}{url}
\documentclass[
]{article}
\usepackage{xcolor}
\usepackage{amsmath,amssymb}
\setcounter{secnumdepth}{-\maxdimen} % remove section numbering
\usepackage{iftex}
\ifPDFTeX
  \usepackage[T1]{fontenc}
  \usepackage[utf8]{inputenc}
  \usepackage{textcomp} % provide euro and other symbols
\else % if luatex or xetex
  \usepackage{unicode-math} % this also loads fontspec
  \defaultfontfeatures{Scale=MatchLowercase}
  \defaultfontfeatures[\rmfamily]{Ligatures=TeX,Scale=1}
\fi
\usepackage{lmodern}
\ifPDFTeX\else
  % xetex/luatex font selection
\fi
% Use upquote if available, for straight quotes in verbatim environments
\IfFileExists{upquote.sty}{\usepackage{upquote}}{}
\IfFileExists{microtype.sty}{% use microtype if available
  \usepackage[]{microtype}
  \UseMicrotypeSet[protrusion]{basicmath} % disable protrusion for tt fonts
}{}
\makeatletter
\@ifundefined{KOMAClassName}{% if non-KOMA class
  \IfFileExists{parskip.sty}{%
    \usepackage{parskip}
  }{% else
    \setlength{\parindent}{0pt}
    \setlength{\parskip}{6pt plus 2pt minus 1pt}}
}{% if KOMA class
  \KOMAoptions{parskip=half}}
\makeatother
\setlength{\emergencystretch}{3em} % prevent overfull lines
\providecommand{\tightlist}{%
  \setlength{\itemsep}{0pt}\setlength{\parskip}{0pt}}
\usepackage{bookmark}
\IfFileExists{xurl.sty}{\usepackage{xurl}}{} % add URL line breaks if available
\urlstyle{same}
\hypersetup{
  pdftitle={CytoReason R to Python},
  hidelinks,
  pdfcreator={LaTeX via pandoc}}

\title{CytoReason R to \texttt{Python}}
\author{}
\date{}

\begin{document}
\maketitle

\subsection{Description}\label{description}

This is a customized syllabus for CytoReason.

This course is intended for R programmers who need to become productive
in \texttt{Python} for data science and bioinformatics work. Rather than
teaching \texttt{Python} from scratch, the course constantly maps
\texttt{Python} concepts back to their R equivalents --- 0-based
indexing versus 1-based, list comprehensions versus \texttt{apply},
dictionaries versus named lists, \texttt{pandas} versus \texttt{dplyr},
and so on --- so that existing R knowledge is leveraged rather than
discarded.

The emphasis of the course is on \texttt{Python}, \texttt{numpy},
\texttt{pandas} and statistical work with \texttt{scikit-learn}. The
course works through a realistic data pipeline: cleaning data, building
\texttt{DataFrame} pipelines, merging clinical metadata with sequence
count matrices, dimensionality reduction with \texttt{PCA}, parsing
memory-heavy biological file formats, optimizing and parallelizing heavy
computation, and finally packaging the whole thing as a robust
command-line tool.

Each module pairs a focused teaching block with a hands-on exercise that
ports or rebuilds a realistic R workflow in \texttt{Python}.

\subsection{Duration}\label{duration}

16 hours / 2 days

\subsection{Intended Audience}\label{intended-audience}

\begin{itemize}
\tightlist
\item
  R programmers moving to \texttt{Python} for data science or
  bioinformatics work.
\item
  Data scientists who need to read, maintain and write \texttt{Python}
  data pipelines.
\end{itemize}

\subsection{Prerequisites}\label{prerequisites}

\begin{itemize}
\tightlist
\item
  Working knowledge of R, including \texttt{dplyr} / tidyverse style
  data manipulation.
\item
  General familiarity with data analysis concepts such as data frames,
  joins and aggregations.
\end{itemize}

\subsection{Objectives}\label{objectives}

\begin{itemize}
\tightlist
\item
  Map core R concepts and idioms onto their \texttt{Python} equivalents.
\item
  Replace tidyverse data pipelines with idiomatic \texttt{pandas}.
\item
  Merge, reshape and clean real-world data sets in \texttt{pandas}.
\item
  Set up reproducible \texttt{Python} environments and model data with
  classes.
\item
  Run \texttt{scikit-learn} pipelines including scaling and
  \texttt{PCA}.
\item
  Parse large biological files efficiently using \texttt{Biopython} and
  generators.
\item
  Optimize and parallelize heavy numerical computation.
\item
  Turn an analysis workflow into a robust command-line tool.
\end{itemize}

\subsection{Outline}\label{outline}

\begin{itemize}
\tightlist
\item
  \texttt{Python} for R users

  \begin{itemize}
  \tightlist
  \item
    0-based indexing and what changes because of it
  \item
    Lists, tuples and slicing versus R vectors
  \item
    List comprehensions versus \texttt{apply} family functions
  \item
    Dictionaries versus named lists
  \item
    Strings, formatting and f-strings
  \item
    Functions, default arguments and keyword arguments
  \item
    Reading and cleaning a basic data set
  \item
    Exercise: porting a basic R data-cleaning script to \texttt{Python}
  \end{itemize}
\end{itemize}

\begin{itemize}
\tightlist
\item
  \texttt{pandas} for \texttt{dplyr} users

  \begin{itemize}
  \tightlist
  \item
    The \texttt{DataFrame} and \texttt{Series} data structures
  \item
    Selecting columns and rows
  \item
    \texttt{loc} versus \texttt{iloc}
  \item
    Filtering rows by condition
  \item
    Adding and transforming columns
  \item
    Method chaining versus the pipe operator
  \item
    Mapping common \texttt{dplyr} verbs onto \texttt{pandas}
  \item
    Exercise: replicating a tidyverse pipeline in \texttt{pandas}
  \end{itemize}
\end{itemize}

\begin{itemize}
\tightlist
\item
  Data merging and reshaping

  \begin{itemize}
  \tightlist
  \item
    Joins with \texttt{merge} and the different join types
  \item
    Concatenating and stacking data
  \item
    Aggregations with \texttt{groupby}
  \item
    Pivoting, melting and reshaping
  \item
    Handling missing data
  \item
    Indexes and multi-indexes
  \item
    Exercise: merging clinical metadata with sequence count matrices
  \end{itemize}
\end{itemize}

\begin{itemize}
\tightlist
\item
  Environments and \texttt{OOP}

  \begin{itemize}
  \tightlist
  \item
    Installing packages with pip and \texttt{conda}
  \item
    Virtual environments and why they matter
  \item
    Reproducible environments and dependency files
  \item
    Classes versus \texttt{S3} / \texttt{S4} objects
  \item
    Attributes, methods and the \texttt{self} parameter
  \item
    Constructors and basic data modeling with classes
  \item
    Exercise: setting up a reproducible environment and writing a custom
    data class
  \end{itemize}
\end{itemize}

\begin{itemize}
\tightlist
\item
  Dimensionality reduction

  \begin{itemize}
  \tightlist
  \item
    The \texttt{scikit-learn} estimator and transformer \texttt{API}
  \item
    Data scaling and standardization
  \item
    \texttt{PCA} concepts and how to apply it
  \item
    Explained variance and choosing components
  \item
    Visualizing results with \texttt{matplotlib}
  \item
    Exercise: running a \texttt{PCA} pipeline on the merged matrices and
    visualizing variance
  \end{itemize}
\end{itemize}

\begin{itemize}
\tightlist
\item
  Specialized bio-tooling

  \begin{itemize}
  \tightlist
  \item
    Overview of \texttt{Biopython}
  \item
    Parsing \texttt{FASTA} and \texttt{VCF} files
  \item
    Iterators and generator expressions
  \item
    Why generators matter for memory-heavy files
  \item
    Streaming records instead of loading everything into memory
  \item
    Exercise: parsing memory-heavy \texttt{FASTA} / \texttt{VCF} files
    efficiently
  \end{itemize}
\end{itemize}

\begin{itemize}
\tightlist
\item
  Code optimization and parallel processing

  \begin{itemize}
  \tightlist
  \item
    \texttt{numpy} arrays and vectorization
  \item
    Replacing loops with vectorized operations
  \item
    Profiling code to find the real bottlenecks
  \item
    Parallelism with \texttt{joblib}
  \item
    The \texttt{multiprocessing} module and the \texttt{GIL}
  \item
    Distributing a heavy calculation across multiple \texttt{CPU} cores
  \item
    Exercise: profiling code and distributing a heavy calculation across
    multiple \texttt{CPU} cores
  \end{itemize}
\end{itemize}

\begin{itemize}
\tightlist
\item
  Pipeline automation

  \begin{itemize}
  \tightlist
  \item
    Parsing command-line arguments with \texttt{argparse}
  \item
    Structuring a script as a reusable tool
  \item
    Error handling and exceptions
  \item
    Logging with the \texttt{logging} module
  \item
    Making a workflow robust and reproducible
  \item
    Exercise: converting the \texttt{PCA} workflow into a robust
    command-line tool (\texttt{CLI})
  \end{itemize}
\end{itemize}

\subsection{Installations}\label{installations}

Each student should have:

\begin{itemize}
\tightlist
\item
  An \texttt{Ubuntu} 24:04 machine, real or virtual.
\item
  4 GB \texttt{RAM} for each machine is enough. This is not much.
\item
  Free, wide band, access to the internet from all machines with no
  weird corporate firewalls that might stop us from installing software
  and \texttt{Python} packages via pip.
\item
  Username and password of a user that has \texttt{sudo} privileges on
  the machine.
\item
  \url{https://www.linuxvmimages.com/images/ubuntu-2204/}
\item
  Users who want to exercise on \texttt{Windows} or on \texttt{MacOS}
  are welcome to do so and I will help them to fix issues but a minority
  of the more advanced exercises may not work on those machine.
\end{itemize}

\subsection{Copyright}\label{copyright}

Mark Veltzer
\href{mailto:mark.veltzer@gmail.com}{\nolinkurl{mark.veltzer@gmail.com}},
© 2026

\end{document}
